\section{Appendix V: Effective Field Theory Interpretation}

This appendix clarifies the domain of validity of the TRXT framework, addressing critiques regarding non-renormalizability of the NJL model at the Planck scale.

\subsection{Statement of EFT Validity}

The Nambu-Jona-Lasinio (NJL) Lagrangian employed in this work is understood as an \textbf{effective field theory (EFT)} valid below a UV cutoff $\Lambda_{UV}$. We explicitly state:

\begin{mdframed}[backgroundcolor=yellow!10]
\textbf{EFT Validity Statement:} The TRXT framework is defined as an effective theory valid for energies $E < \Lambda_{UV}$, where:
\begin{equation}
    \Lambda_{UV} \sim 0.1 \, M_{Pl} \approx 10^{18} \text{ GeV}
\end{equation}
Above this scale, the theory requires UV completion (e.g., via Asymptotic Safety, String Theory, or Loop Quantum Gravity).
\end{mdframed}

\subsection{Power Counting and Operator Classification}

The NJL action with contact four-fermion interaction:
\begin{equation}
    S_{NJL} = \int d^4x \left[ \bar{\psi} i \slashed{\partial} \psi + \frac{G}{2} (\bar{\psi}\psi)^2 \right]
\end{equation}
has coupling dimension $[G] = -2$ (mass dimension), making it \textbf{power-counting non-renormalizable}. However, within EFT logic:

\begin{enumerate}
    \item \textbf{Relevant Operators ($[\mathcal{O}] < 4$):} Mass terms $m\bar{\psi}\psi$ — generated dynamically via chiral symmetry breaking.
    \item \textbf{Marginal Operators ($[\mathcal{O}] = 4$):} Kinetic term $\bar{\psi} i \slashed{\partial} \psi$ — renormalizable.
    \item \textbf{Irrelevant Operators ($[\mathcal{O}] > 4$):} Four-fermion $G(\bar{\psi}\psi)^2$ — suppressed by $E^2/\Lambda_{UV}^2$ at low energies.
\end{enumerate}

At energies $E \ll \Lambda_{UV}$, irrelevant operators contribute corrections of order $(E/\Lambda_{UV})^{n-4}$, which are negligible for $n > 4$.

\subsection{Asymptotic Safety Perspective}

An alternative UV completion, consistent with TRXT philosophy, is Weinberg's Asymptotic Safety program \cite{weinberg_as}:

\textbf{Hypothesis:} There exists a non-trivial UV fixed point where:
\begin{equation}
    G(\mu) \to G^* \quad \text{as} \quad \mu \to \infty
\end{equation}
At the fixed point, the theory is scale-invariant and well-defined to arbitrary energies.

Evidence from functional renormalization group (FRG) studies suggests that gravity + matter systems can exhibit such fixed points \cite{reuter}.

\subsection{Implications for TRXT Predictions}

\begin{enumerate}
    \item \textbf{Low-Energy Predictions (Valid):} Galaxy rotation curves, SIDM cross-sections, neutrino phenomenology — all lie within EFT regime and are robust.
    \item \textbf{Planck-Scale Physics (Speculative):} Big Condensation scenario, exact vacuum energy calculation — require UV completion or Asymptotic Safety.
    \item \textbf{Error Estimate:} For observables at $E \sim 1$ GeV, EFT corrections are $\sim (1/10^{18})^2 \approx 10^{-36}$ — utterly negligible.
\end{enumerate}

\subsection{Conclusion}

The TRXT framework is self-consistent as an EFT below $\Lambda_{UV} \sim 0.1 M_{Pl}$. Claims about Planck-scale physics (Layers 0–2 in Section 2) are understood as \textit{theoretical motivation}, not rigorous predictions. All falsifiable predictions (SPARC, SIDM, MaVaN) lie well within the EFT validity regime.
